\chapter*{Preface}

Data arrive as tables of numbers, and almost everything we do with them is linear algebra, optimisation or
probability. This book proves the theorems behind that work: why the algorithms that compute with matrices can be
trusted, why the methods that minimise functions converge and how fast they can possibly converge, what a sample of
data can and cannot tell us, and why a machine that has seen finitely many examples can predict the next one.

The first two chapters compute. Rounding errors are tamed by backward error analysis, which explains why Householder
reflections, the $QR$ algorithm and the singular value decomposition are reliable, and Krylov spaces and random
projections make enormous problems small. The next four optimise. Convex duality turns every minimisation into a
maximisation that certifies it; gradient descent, Nesterov's acceleration and their lower bounds show exactly how fast
first-order methods can go; proximal methods find sparse solutions, and with them compressed sensing recovers a signal
from far fewer measurements than its length; Newton's method and interior points solve linear programs in polynomial
time. The middle of the book learns from data that arrive one at a time, against an adversary, and then turns to
statistics: estimation, the Cram\'er--Rao bound and the startling paradox of James and Stein, tests and confidence
intervals, the false discovery rate and the bootstrap, and Bayesian inference computed by Markov chain Monte Carlo.
The last two chapters explain generalisation, through the Vapnik--Chervonenkis dimension and boosting, and the two great
families of learning machines, kernels and neural networks, up to Barron's theorem that a two-layer network escapes the
curse of dimension.

The prerequisites are Tomes~0, 3, 6 and~7, with Tome~S1 for linear algebra and Tome~C2 for linear programming; results
from them are recalled where they are used, with a precise reference. Every theorem is proved, including the error
bounds and rates of convergence, which are theorems like any other. Each chapter opens with a story and a question to
think about, and ``Pause and think'' boxes ask you to predict before reading on. Exercises are graded from $\star$ to
$\star\star\star$; each has a hint and a solution that begins by explaining how one might find the idea.

Each chapter has a companion Python notebook that implements its methods from scratch and checks every number in the
text, so that the reader can watch rounding errors stay below their bounds, a zigzagging gradient method lose its race
against an accelerated one, a sparse signal reappear from a handful of random measurements, and a confidence interval
miss the truth exactly as often as it should.
